%% Separately authored consumer-corrected historical derivative.
%% Necessity: the exact source restores the full CH/non-CH results but does not
%% contain the broader Boolean-valued, permutation-model, and independence catalog.
%% Definition: this file is based on historical forcing-variants lines 1--268;
%% source consumers are redirected, captions omit internal IDs, and the
%% Boolean-existential sethood bridge is supplied explicitly.
%% Standing: this is separately authored reader material, not inherited source
%% or a verification of the imported forcing theorem.
%% Nonredundancy: the exact source remains controlling for CH/non-CH; this file
%% supplies broader variant definitions and a cross-topic catalog.
%% Use: parts/p06.tex routes this file once after the exact source and bounded
%% forcing apparatus, and route tests pin the complete corrected derivative.

% BEGIN HISTORICAL SLICE 1-268
%% parts/p06/forcing-variants.tex --- P6 forcing catalogs and variants

\begingroup
\pagestyle{plain}% Suppress stale section marks on both A-067 landscape pages.
\begin{landscape}
\thispagestyle{plain}% Keep the rotated running head clear of the wide A-067 table.
\section{Forcing catalogs and symmetric models}
\label{sec:p06-forcing-variants}

\noindent\textit{Reader note.}  This separately authored comparison section
supplements the exact CH/non-CH forcing route with Boolean-valued and symmetric
model presentations.  Its cited variants do not replace the exact forcing
definitions or their proof boundaries.

The countable-transitive-model presentation in
Chapter~\ref{ch:constructibility-forcing} is the pedagogical route used
for the migrated source proofs.  This section supplies the comparison entries
that those proofs do not: exact poset conventions, the Boolean-valued
presentation, and the choice-failure method.  Each independence claim below is
relative to the stated background theory.

\subsection{Three forcing conventions}
\scriptsize
\setlength{\tabcolsep}{2.5pt}
\begin{longtable}{@{}p{.58in}p{1.78in}p{.58in}p{.68in}p{1.00in}
  p{1.00in}p{1.42in}@{}}
\caption{Forcing-poset catalog.  The order always points
toward stronger information; collapse notation is fixed in the row rather
than presumed universal.}\label{tab:p06-forcing-catalog}\\
\toprule
Forcing & Exact poset & Chain condition & Closure & Adds & Collapses & Preserves \\
\midrule
\endfirsthead
\toprule
Forcing & Exact poset & Chain condition & Closure & Adds & Collapses & Preserves \\
\midrule
\endhead
Cohen
& For an infinite cardinal \(\kappa\),
  \(\operatorname{Add}(\omega,\kappa)
    =\operatorname{Fn}(\kappa\times\omega,2,{<}\omega)\), ordered by reverse
  inclusion.
& ccc
& Not countably closed
& \(\kappa\) distinct coordinate Cohen reals; if
  \(\kappa^{\aleph_0}=\kappa\), nice-name counting gives
  \(2^{\aleph_0}=\kappa\).
& None
& All ground-model cardinals and cofinalities, by ccc. \\
\addlinespace
Random
& The complete measure algebra: the sigma-algebra generated by the basic
  clopen cylinders in \(2^\omega\), modulo the ideal of sets of product
  measure zero; equivalently, use positive Borel representatives ordered by
  inclusion modulo null sets.
& ccc
& Not countably closed
& One random real; Boolean values are null-equivalence classes.
& None
& All ground-model cardinals and cofinalities, by ccc. \\
\addlinespace
Collapse
& For regular cardinals \(\kappa<\lambda\),
  \(\operatorname{Coll}(\kappa,\lambda)\defeq\) the set of partial functions
  \(p:\kappa\rightharpoonup\lambda\) of size \(<\kappa\), ordered by reverse
  inclusion.
& Size \(\lambda^{<\kappa}\), hence
  \((\lambda^{<\kappa})^+\)-cc; if \(\lambda^{<\kappa}=\lambda\), then
  \(\lambda^+\)-cc.
& \({<}\kappa\)-closed
& A generic surjection \(\kappa\twoheadrightarrow\lambda\)
& Every ground cardinal in \((\kappa,\lambda]\) to size \(\kappa\);
  \(\kappa\) itself is preserved
& Cardinals at most \(\kappa\); if \(\lambda^{<\kappa}=\lambda\), cardinals
  above \(\lambda\).  No blanket cofinality claim. \\
\bottomrule
\end{longtable}
\end{landscape}
\endgroup

The claims in Table~\ref{tab:p06-forcing-catalog} are cited to
\textcite[pp.~201--265]{Jech2003} and
\textcite[chs.~2, 4, and 5]{Bell2005}.  In particular, the collapse row does
not make an unqualified preservation claim.

\subsection{Boolean-valued presentation}

\begin{hbdefinition}{def:p06-boolean-valued-universe}{Boolean-valued universe}
Let \(\mathbb B\) be a complete Boolean algebra, meaning that every subset of
\(\mathbb B\) has a join and a meet.  Write
\(b^{\complement_{\mathbb B}}\), \(b\wedge_{\mathbb B}c\), and
\(b\vee_{\mathbb B}c\) for Boolean complement, meet, and join; these
subscripts distinguish Boolean-algebra operations from the globally
registered object-language connectives.  Define recursively
\[
  V^{\mathbb B}_\alpha
  \defeq\left\{\mathbf u:
      \begin{array}{l}
      \mathbf u\text{ is a function},\quad
      \operatorname{ran}(\mathbf u)\subseteq\mathbb B,\\
      \dom(\mathbf u)\subseteq
        \displaystyle\bigcup_{\beta<\alpha}V^{\mathbb B}_\beta
      \end{array}\right\}.
\]
Put
\[
  V^{\mathbb B}\defeq
  \bigcup_{\alpha\in\operatorname{Ord}}V^{\mathbb B}_\alpha.
\]
The members of \(V^{\mathbb B}\) are \(\mathbb B\)-names.  The Boolean truth
map \(\operatorname{Bool}_{\mathbb B}\) is the simultaneous recursion
\[
\begin{aligned}
 \operatorname{Bool}_{\mathbb B}(\mathbf u\in\mathbf v)
   &\defeq\mathop{\bigvee}\nolimits^{\mathbb B}_{\mathbf w\in\dom(\mathbf v)}
       \mathbf v(\mathbf w)\wedge_{\mathbb B}
       \operatorname{Bool}_{\mathbb B}(\mathbf u=\mathbf w),\\
 \operatorname{Bool}_{\mathbb B}(\mathbf u=\mathbf v)
   &\defeq\mathop{\bigwedge}\nolimits^{\mathbb B}_{\mathbf w\in\dom(\mathbf u)}
       \bigl(\mathbf u(\mathbf w)^{\complement_{\mathbb B}}\vee_{\mathbb B}
             \operatorname{Bool}_{\mathbb B}(\mathbf w\in\mathbf v)\bigr)\\
   &\quad\wedge_{\mathbb B}
     \mathop{\bigwedge}\nolimits^{\mathbb B}_{\mathbf w\in\dom(\mathbf v)}
       \bigl(\mathbf v(\mathbf w)^{\complement_{\mathbb B}}\vee_{\mathbb B}
             \operatorname{Bool}_{\mathbb B}(\mathbf w\in\mathbf u)\bigr),\\
 \operatorname{Bool}_{\mathbb B}(\neg\varphi)
   &\defeq
     \operatorname{Bool}_{\mathbb B}(\varphi)^{\complement_{\mathbb B}},\\
 \operatorname{Bool}_{\mathbb B}(\varphi\wedge\psi)
   &\defeq\operatorname{Bool}_{\mathbb B}(\varphi)\wedge_{\mathbb B}
       \operatorname{Bool}_{\mathbb B}(\psi),\\
 \operatorname{Bool}_{\mathbb B}(\exists x\,\varphi(x,\bar{\mathbf v}))
   &\defeq\mathop{\bigvee}\nolimits^{\mathbb B}
       S_{\varphi,\bar{\mathbf v}}.
\end{aligned}
\]
At the existential step of the simultaneous recursion, define
\[
 S_{\varphi,\bar{\mathbf v}}
 \defeq
 \left\{b\in\mathbb B:
   \exists\mathbf u\in V^{\mathbb B}\,
   b=\operatorname{Bool}_{\mathbb B}
      (\varphi(\mathbf u,\bar{\mathbf v}))\right\}.
\]
This is a set by Separation because it is a subclass of the set
\(\mathbb B\); completeness of \(\mathbb B\) therefore supplies the displayed
join.  Equivalently one may rank-bound the witness names and prove that the
value stabilizes above the parameter ranks.  The class
\(V^{\mathbb B}\) is never used as a primitive proper-class index for a
Boolean join \cite{ScottSolovay1967,Bell2005}.

The remaining connectives and \(\forall\) are defined from these.  Completeness
of \(\mathbb B\) is needed for the displayed set-indexed joins and meets.  This map has
Boolean-algebra values and is neither the two-valued satisfaction relation nor
the forcing relation of Chapter~\ref{ch:constructibility-forcing}.  Unlike
a \(\mathbb P\)-name, a \(\mathbb B\)-name is a Boolean-valued function on
earlier names, not a set of name--condition pairs; this distinction is exactly
what the generic-quotient theorem uses.
\index{Boolean-valued model!truth map}
\end{hbdefinition}

\begin{hbconvention}{def:p06-internal-boolean-universe}{Boolean-valued recursion computed in a ground model}
Let \(M\) be transitive, let \(\mathbb B\in M\), and suppose
\(M\entails\) “\(\mathbb B\) is a complete Boolean algebra.”  Write
\((V^{\mathbb B})^M\) for the name hierarchy constructed inside \(M\) and
\(\operatorname{Bool}_{\mathbb B}^M\) for the simultaneous truth recursion
computed there.  In its existential clause,
\(S_{\varphi,\bar{\mathbf v}}^M\) is the subset of \(\mathbb B\) obtained by
Separation in \(M\), and \(\bigvee^M S^M\) is the join supplied by the
completeness that \(M\) computes.

The superscript is necessary because an algebra complete in \(M\) need not be
complete in the ambient universe.  The notation stands for internal classes
and operations of \(M\), differs from the ambient-complete construction in
Definition~\ref{def:p06-boolean-valued-universe}, and is used by the quotient
theorem below.
\end{hbconvention}

\begin{hbtheorem}{thm:p06-boolean-valued-quotient}{cited=Bell2005}{Generic quotient of a Boolean-valued universe}
Let \(M\) be a transitive model of \(\mathsf{ZFC}\), let
\(\mathbb B\in M\) satisfy \(M\entails\) “\(\mathbb B\) is complete,” and
put \(\mathbb B^+=\mathbb B\setminus\{0_{\mathbb B}\}\), ordered by the
Boolean order so that smaller nonzero elements are stronger conditions.
Let \(G\subseteq\mathbb B^+\) be a poset filter meeting every dense
\(D\in M\) with \(D\subseteq\mathbb B^+\).  Equivalently, under the
Boolean-algebra reading this same \(G\) is an \(M\)-generic ultrafilter on
\(\mathbb B\), with \(0_{\mathbb B}\notin G\).

Use \((V^{\mathbb B})^M\) and
\(\operatorname{Bool}_{\mathbb B}^M\) from
Convention~\ref{def:p06-internal-boolean-universe}.  Quotienting those
\(M\)-names by
\[
  \mathbf u\sim_G\mathbf v
  \quad\defeq\quad
  \operatorname{Bool}_{\mathbb B}^M(\mathbf u=\mathbf v)\in G
\]
and interpreting membership by
\(\operatorname{Bool}_{\mathbb B}^M(\mathbf u\in\mathbf v)\in G\) yields the ordinary
two-valued forcing extension \(M[G]\).  The cited forcing theorem transfers
\(\mathsf{ZFC}\) and preserves the ordinals.
\end{hbtheorem}

Products coordinate several forcing notions at once; finite-support iterations
feed each successive extension into the next stage.  These are navigation
pointers only.  Proper forcing, countable-support iterations, the Proper
Forcing Axiom \(\mathsf{PFA}\), and Martin's Maximum \(\mathsf{MM}\) are
outside the scope of this forcing survey.

\subsection{Permutation and symmetric models}

\begin{hbdefinition}{def:p06-permutation-model-datum}{Permutation data and support}
A permutation-model datum is a triple \((A,\mathcal G,\mathfrak F)\), where
\(A\) is a set of atoms, \(\mathcal G\leq\operatorname{Sym}(A)\), and
\(\mathfrak F\) is a normal filter of subgroups of \(\mathcal G\): it is a
nonempty family, upward closed under supergroups within \(\mathcal G\), closed
under finite intersections, and closed under conjugation by members of
\(\mathcal G\); moreover, for every \(a\in A\), it contains the pointwise
  stabilizer \(\{\pi\in\mathcal G:\pi(a)=a\}\).  Extend the action to the
  cumulative hierarchy over \(A\) by
\(\pi x=\{\pi y:y\in x\}\) for sets, retaining the given action on atoms.  Put
\(\operatorname{Stab}_{\mathcal G}(x)
 \defeq\{\pi\in\mathcal G:\pi x=x\}\).  A set \(x\) is
\emph{symmetric} iff \(\operatorname{Stab}_{\mathcal G}(x)\in\mathfrak F\), and is
\emph{hereditarily symmetric} iff \(x\) and every member of its transitive
closure are symmetric.  For the finite-support filter, a finite
\(E\subseteq A\) supports
\(x\) iff the pointwise stabilizer of \(E\) is contained in the stabilizer of
\(x\).

The group and filter are required to distinguish the retained sets from the
full ZFA universe; normality makes the retained class invariant, and the
hereditary clause makes it transitive.  These distinctions are used in
Theorem~\ref{thm:p06-basic-fraenkel-model}.
\index{permutation model!support}
\end{hbdefinition}

\begin{hbtheorem}{thm:p06-basic-fraenkel-model}{cited=Halbeisen2017}{The basic Fraenkel model of ZFA}
Externally let \(A\) be a countably infinite set of atoms, take
\(\mathcal G=\operatorname{Sym}(A)\), and let \(\mathfrak F\) be generated by
the pointwise stabilizers of finite subsets of \(A\).  In the resulting
hereditarily symmetric model of \(\mathsf{ZFA}\), every subset of \(A\) is
finite or cofinite.  Hence \(A\) is amorphous---infinite, but with every subset
finite or cofinite---and cannot be well ordered; therefore Choice fails.
\end{hbtheorem}

The atom model is not itself a model of pure \(\mathsf{ZF}\).  Transfer from
permutation models to pure \(\mathsf{ZF}\), or a direct symmetric-forcing
construction, is a cited step; see \textcite{Jech1973Choice,Halbeisen2017}.
The historical permutation construction is due to \textcite{Fraenkel1922}.

\begin{landscape}
\subsection{Relative-independence catalog}
\small
\begin{longtable}{@{}p{.85in}p{1.95in}p{2.10in}p{2.10in}@{}}
\caption{Relative-independence catalog.  ``Positive'' and
``negative'' name relative-consistency directions, not truth values in
\(\mathsf{ZFC}\).}\label{tab:p06-independence-catalog}\\
\toprule
Statement & Positive direction & Negative direction & Method and evidence \\
\midrule
\endfirsthead
\toprule
Statement & Positive direction & Negative direction & Method and evidence \\
\midrule
\endhead
\(\mathsf{CH}\)
& \(\mathsf{Con}(\mathsf{ZFC})\) implies
  \(\mathsf{Con}(\mathsf{ZFC}+\mathsf{CH})\).
& \(\mathsf{Con}(\mathsf{ZFC})\) implies
  \(\mathsf{Con}(\mathsf{ZFC}+\neg\mathsf{CH})\).
& Positive side: the local \(L\)-relativization argument and
  \textcite{Godel1938}.  Negative side: the cited Boolean-valued forcing
  import \parencite{Cohen1963,Jech2003}.  The final nonderivability step is
  Theorem~\ref{thm:ch-nonderivability}. \\
\addlinespace
\(\mathsf{GCH}\) patterns
& \(\mathsf{Con}(\mathsf{ZFC})\) implies
  \(\mathsf{Con}(\mathsf{ZFC}+\mathsf{GCH})\), by passage to \(L\).
& \(\mathsf{Con}(\mathsf{ZFC})\) implies
  \(\mathsf{Con}(\mathsf{ZFC}+\neg\mathsf{GCH})\).  More generally, over a
  \(\mathsf{ZFC}+\mathsf{GCH}\) ground model, Easton forcing realizes each
  definable regular-cardinal function \(F\) satisfying monotonicity and
  \(\operatorname{cf}(F(\kappa))>\kappa\); for example,
  \(F(\kappa)=\kappa^{++}\).
& Positive side: constructibility and \textcite{Godel1938}.  Negative side:
  \textcite{Easton1970}.  No row claims arbitrary freedom at singular
  cardinals. \\
\addlinespace
\hyperref[def:p06-suslin-hypothesis]{Chapter~6.4 tree assertion}
& \(\mathsf{Con}(\mathsf{ZFC})\) implies consistency of \(\mathsf{ZFC}\) with
  the nonexistence of an \(\omega_1\)-tree having countable levels and neither
  an uncountable chain nor an uncountable antichain, via a cited forcing model.
& \(\mathsf{Con}(\mathsf{ZFC})\) implies consistency of \(\mathsf{ZFC}\) with
  the existence of such a tree, via a cited constructible-universe model.
& Positive side: \textcite{MartinSolovay1970,Jech2003}.  Negative side:
  \textcite{Jensen1972FineStructure}.  The required iteration and inner-model
  arguments are not reproduced here; Chapter~6.4 supplies the standard names
  and exact definitions. \\
\addlinespace
Choice over \(\mathsf{ZF}\)
& \(\mathsf{Con}(\mathsf{ZF})\) implies
  \(\mathsf{Con}(\mathsf{ZF}+\mathsf{AC})\), by passage to the
  constructible inner model.
& \(\mathsf{Con}(\mathsf{ZF})\) implies
  \(\mathsf{Con}(\mathsf{ZF}+\neg\mathsf{AC})\), by permutation or
  symmetric models together with the cited transfer to pure set theory.
& Positive side: constructibility and \textcite{Godel1938}.  Negative side:
  \textcite{Jech1973Choice,Halbeisen2017}.  The local Fraenkel example is only
  a \(\mathsf{ZFA}\) permutation model; passage to pure \(\mathsf{ZF}\) uses
  the cited atoms-to-pure-set transfer. \\
\bottomrule
\end{longtable}
\end{landscape}
% END HISTORICAL SLICE 1-268
